% !TEX root = ../thesis_book_0421052099.tex
\chapter{Tool API Specification and Cypher Templates}\label{app:toolapi}

% Cypher reproduced from agr/kg_tools.py. The parameterisation is the point: the
% model supplies only $-parameters, never query text.

\Cref{sec:tool-api} argues that the model must never author a query. This
appendix gives the five operations that follow from that decision, with the
Cypher each one executes. Every query below is a fixed template. The model
supplies values for the \verb|$|-parameters and nothing else, so a malformed
model output produces a failed lookup rather than a syntax error, and the error
taxonomy of \Cref{ch:erroranalysis} never has to separate reasoning failures
from query-authoring failures.

Two limits appear throughout. \textsf{max\_relations}, set to $300$, caps the
relation list returned for an entity; \textsf{max\_fanout}, set to $200$, caps
the neighbours returned for one edge. \Cref{tab:rootcauses} records why the
first was raised from its original value of $80$.

Every invocation is appended to a JSONL log with its arguments, its result, and
its wall-clock duration in milliseconds. That log is the per-call record the
latency figures of \Cref{sec:efficiency} are computed from, and for the
determinism probes it is the only surviving per-question artifact.

\section{search\_entity}\label{app:tool-search}

Resolves a surface form to graph entities. This is the only entry point from
text to graph, and the only operation that consults an embedding. It delegates
to the three-stage resolver of \Cref{sec:entity-resolution} --- exact, lexical,
vector --- and returns which stage produced the hit, so that a downstream
failure can be attributed to resolution rather than to navigation.

One naming collision has to be declared rather than tidied away. The returned
JSON key is literally \textsf{tier}, and it takes the values $1$, $2$ and $3$
for the three stages. \Cref{sec:groundedness-metric} reserves that word for the
groundedness metric, and this is the sole exception: the key appears verbatim in
every committed tool log and in the scripts that read them, so renaming it in
this appendix would misdocument the artifacts. The two senses are unrelated ---
a resolver stage is not a groundedness tier --- and the stage names are used
everywhere the prose has a choice.

\begin{lstlisting}[style=verbatimblock]
search_entity(surface_form, k=5)
  -> [{id, name, score, tier}, ...]     up to k candidates

tier=1  exact:     MATCH (e:Entity {name:$m}) RETURN e.id, e.name
tier=2  lexical:   CALL db.index.fulltext.queryNodes('entity_name', $q)
                     YIELD node, score
                   RETURN node.id, node.name, score LIMIT $k
tier=3  vector:    CALL db.index.vector.queryNodes('entity_embedding', $k, $v)
                     YIELD node, score
                   RETURN node.id, node.name, score
\end{lstlisting}

The lexical stage accepts its top hit only when the full-text score exceeds a
threshold of $1.0$. Without that gate the index returns a ranked list for almost
any input, including inputs with no correct answer in the graph, and the vector
stage would never be reached.

\section{get\_relations}\label{app:tool-relations}

Returns the relation types incident to an entity, in \emph{both} directions,
each with its fanout count. Direction is carried explicitly because Freebase
edges are asserted in one direction only, and the answer to a question is as
often reached by traversing an edge backwards as forwards.

\begin{lstlisting}[style=verbatimblock]
get_relations(entity_id)
  -> [{rel, dir, n}, ...]     sorted by n descending, capped at max_relations

MATCH (e:Entity {id:$id})-[r]->()
RETURN r.fb_name AS rel, 'out' AS dir, count(*) AS n
UNION ALL
MATCH (e:Entity {id:$id})<-[r]-()
RETURN r.fb_name AS rel, 'in' AS dir, count(*) AS n
\end{lstlisting}

Administrative predicates are filtered before the cap is applied, by prefix:
\textsf{common.topic.}, \textsf{common.}, \textsf{freebase.}, \textsf{type.},
\textsf{kg.}, \textsf{user.}, \textsf{dataworld.}, \textsf{rdf-schema\#} and
\textsf{owl\#}. These are schema bookkeeping rather than facts about the world.
The list is quoted as the code writes it, which is why it holds nine entries
where \Cref{sec:tool-api} names eight: \textsf{common.topic.} is subsumed by
\textsf{common.} and filters nothing the shorter prefix does not already filter.
They also have very high fanout, so filtering before sorting is what keeps them
from consuming the whole list.

\section{get\_neighbors}\label{app:tool-neighbors}

Expands one edge from one entity. This is the operation that makes Freebase's
mediator nodes invisible to the agent: where a neighbour is a mediator, the tool
hops through it in the same call and returns the entity on the far side,
carrying the \textsf{via} chain that reached it.

\begin{lstlisting}[style=verbatimblock]
get_neighbors(entity_id, relation, direction='out')
  -> {neighbors: [{id, name, via[, cvt]}, ...], truncated: bool}

MATCH (e:Entity {id:$id})
MATCH (e)-[r {fb_name:$rel}]->(n)        -- or <-[r]- when direction='in'
RETURN n.id, n.name, n.is_cvt LIMIT $cap

-- for each neighbour with is_cvt = true, one further hop:
MATCH (c:Entity {id:$cid})-[r2]-(m:Entity)
WHERE m.id <> $orig AND m.is_cvt = false
RETURN m.id, m.name, r2.fb_name AS rel2 LIMIT $cap
\end{lstlisting}

The \textsf{via} chain is what makes the traversal auditable. A two-element
\textsf{via} means the result was reached through a mediator, and the
\textsf{cvt} field names the mediator that was crossed.
\Cref{sec:graph-construction} states the two conventions this creates: a
mediator counts as two edges wherever the thesis measures distance in the graph
(the reachability gate, the hop strata), and as one step against the agent's
depth budget, because the agent should not be charged twice for an artifact of
the encoding.

The \textsf{truncated} flag was intended to report whether the fanout cap was
hit, so that a truncated expansion --- a plausible cause of a later failure ---
could be distinguished from a relation that was simply never explored. It does
not deliver that, for the two reasons \Cref{sec:budgets} records: it is returned
to the caller but never written to the tool log, and it is computed on the row
count before mediator expansion while the list is truncated after it. Nothing in
this thesis reads it.

\section{verify\_triple}\label{app:tool-verify-triple}

Asks whether a specific relation holds between two specific entities. It was
built as the relation-aware structural route of the verification layer, and it
is \emph{not called by any node in the final design}: free-text claim relations
do not map onto Freebase predicates, so requiring the relation to match rejected
true claims, and \textsf{verify\_connection} replaced it.
\Cref{sec:five-operations} gives that history and \Cref{sec:structural-check}
gives the two routes the verifier ended up with. It is reproduced here because
it remains part of the tool API, is exercised by the integration tests, and is
the natural primitive for the relation-aware check \Cref{sec:future} proposes.

\begin{lstlisting}[style=verbatimblock]
verify_triple(head_id, relation, tail_id)
  -> {direct, via_cvt, supported}

MATCH (h:Entity {id:$h}), (t:Entity {id:$t})
RETURN
  EXISTS { MATCH (h)-[{fb_name:$rel}]-(t) } AS direct,
  EXISTS { MATCH (h)-[{fb_name:$rel}]-(c:Entity {is_cvt:true})-[]-(t) }
    AS via_cvt
\end{lstlisting}

Both patterns are undirected. A claim is supported if either holds. The
\textsf{via\_cvt} arm exists for the same reason as the mediator hop in
\textsf{get\_neighbors}: a fact stated through a mediator is still a fact, and a
verifier that missed it would reject true claims about exactly the compositional
questions the thesis is about.

\section{verify\_connection}\label{app:tool-verify-connection}

Asks only whether two entities are adjacent, directly or through one mediator.
The relation is not named.

\begin{lstlisting}[style=verbatimblock]
verify_connection(a_id, b_id)
  -> {direct, via_cvt, connected}

MATCH (a:Entity {id:$a}), (b:Entity {id:$b})
RETURN EXISTS { MATCH (a)-[]-(b) } AS direct,
       EXISTS { MATCH (a)-[]-(c:Entity {is_cvt:true})-[]-(b) } AS via_cvt
\end{lstlisting}

This is the fifth operation, and the one not in the original design.
\Cref{sec:tool-api} records that it was added when the verification layer of
\Cref{sec:verification-layer} needed a relation-agnostic adjacency test. Its
permissiveness is deliberate and is also its weakness:
\Cref{sec:verify-failure-modes} analyses wrongful acceptance as a direct
consequence of asking whether two entities are connected at all rather than
whether they are connected in the way the claim asserts.

\endinput
